\documentclass[twocolumn]{aastex631}

\usepackage{amsmath}
\usepackage{multirow}
\usepackage{natbib}
\usepackage{graphicx} 
\usepackage{aas_macros}

\begin{document}

In this work, we addressed the fundamental challenge of consistently quantizing Degenerate Higher-Order Scalar-Tensor (DHOST) theories, whose path to a quantum description is obstructed by the presence of field-dependent gauge symmetries. We posited that these symmetries do not form a closed Lie algebra but rather an open one, which closes only upon use of the classical equations of motion. This structure invalidates standard quantization techniques and necessitates the use of the more powerful Batalin-Vilkovisky (BV) formalism. Our investigation sought to first rigorously establish this algebraic structure and then leverage the corresponding BV framework to constrain the theory at the quantum level.

Our methodology proceeded in a sequence of three logical steps. First, we performed a direct, first-principles calculation of the commutator of two infinitesimal gauge transformations acting on the scalar and metric fields. This allowed us to explicitly verify the open nature of the gauge algebra and to precisely determine its field-dependent structure functions and the coefficients of the terms proportional to the equations of motion. With this information, we then systematically constructed the complete BV-extended action, enlarging the field space to include ghosts and antifields whose interactions are dictated by the open algebra. This construction yielded a classical action that satisfies the cornerstone of the formalism, the classical master equation. Finally, we elevated this to the quantum level, employing the one-loop Quantum Master Equation as a stringent consistency condition. We proposed a general ansatz for one-loop counter-terms built from curvature-squared invariants and used the master equation as an algebraic filter to determine which forms are compatible with the quantum preservation of the gauge symmetry.

Our analysis yielded a series of powerful and predictive results. We first provided definitive proof that the gauge algebra is indeed open, explicitly deriving the structure functions that govern its on-shell closure. This foundational result mandated the successful construction of the BV-extended action, which correctly encodes the intricate gauge structure. The most striking results emerged from the application of the Quantum Master Equation. This consistency requirement was found to act as a powerful selection principle on potential one-loop corrections. We demonstrated that counter-terms constructed from the square of the Weyl tensor are strictly forbidden, leading to a non-renormalization theorem for this class of operators. In contrast, a counter-term involving the Gauss-Bonnet invariant was not only permitted but was uniquely selected. Most remarkably, quantum consistency fixed the functional form of its coupling to the scalar sector to be a specific logarithmic function, $f_G(\phi, X) \propto \log(c_0 - c_1 X)$.

From these results, we have learned that the apparent complexity of an open gauge algebra is not an obstacle to quantization but rather a source of profound physical constraint. This work establishes a direct and predictive link between a theory's classical algebraic structure and the definitive form of its quantum corrections. The parameters $c_0$ and $c_1$, which define the infinitesimal classical gauge transformations, are precisely the same parameters that dictate the functional form of the leading quantum correction. The quantum theory thus retains a precise memory of its classical symmetry. This demonstrates that the Batalin-Vilkovisky formalism is not merely a technical tool but a crucial framework that uncovers the deep consistency conditions governing the theory, transforming the challenge of field-dependent symmetries into a powerful guiding principle for determining the structure of the quantum effective action.

\end{document}
                